\section{System Model and Problem Formulation}
\label{sec:model}
\begin{table}[!t]
\caption{Principal notation}
\label{tab:notation}
\centering
\footnotesize
\begin{tabular}{@{}p{0.23\columnwidth}p{0.70\columnwidth}@{}}
\toprule
Symbol & Meaning \\
\midrule
$G=(\taskset,\edgeset)$ & Workflow DAG and dependency edges \\
$\nodeset$, $g_k$ & Resource nodes and node $k$'s fault domain \\
$w_i,\sigma_i,\mu_i$ & Task workload, CPU demand, and memory demand \\
$c_k,u_k,q_k$ & Compute speed, background load, and recovery queue estimate \\
$x_{ik},y_i$ & Copy-placement indicator and active-replication mode \\
$\copies{i}$ & Nodes hosting copies of task $i$ \\
$p_{ik},\tau_{k\ell}(d)$ & Execution time and transfer delay \\
$D,R_{\mathrm{SLA}}$ & Deadline and declared recovery-delay limit \\
$\cbase,C^{\omega}$ & Nominal and scenario-specific makespan \\
$\rmax,\reff$ & Largest additional delay and effective recovery budget \\
$h_i,s_{ik}$ & Structural task criticality and learned placement score \\
$\rhat,V$ & Recovery surrogate and estimated violation \\
\bottomrule
\end{tabular}
\end{table}

\subsection{Workflow and Computing Continuum}
A workflow is a DAG $G=(\taskset,\edgeset)$. Task $i$ has workload $w_i$,
CPU demand $\sigma_i$, and memory demand $\mu_i$. Dependency $(i,j)$ carries
data volume $d_{ij}$. The resource set $\nodeset$ is partitioned into cloud,
edge, and device tiers. Node $k$ has compute speed $c_k$, CPU capacity $C_k$,
memory capacity $M_k$, background utilization $u_k$, unit compute price
$\rho_k$, failure-rate parameter $\lambda_k$, and fault-domain label $g_k$.
Its residual deployment capacities are
\begin{equation}
 \overline C_k=C_k(1-u_k),\qquad
 \overline M_k=M_k(1-u_k).
\end{equation}
Execution time is deterministic: $p_{ik}=w_i/c_k$. For different nodes,
\begin{equation}
 \tau_{k\ell}(d)=L_{k\ell}+d/B_{k\ell},\qquad k\ne\ell,
\end{equation}
where $L_{k\ell}$ and $B_{k\ell}$ are link latency and bandwidth. A small
configured local-access delay is used when $k=\ell$. The continuum is fixed
for a scheduling instance; elastic node provisioning, energy consumption,
and stochastic execution-time fluctuations are outside this model.

\subsection{Placement and Protection Modes}
Let $x_{ik}\in\{0,1\}$ indicate whether task $i$ has a copy on node $k$, and
let $y_i\in\{0,1\}$ denote its protection mode. Passive tasks have one
deployed copy and can be reexecuted after failure. Active tasks have multiple
concurrently provisioned copies. Define
$\copies{i}=\{k:x_{ik}=1\}$ and impose
\begin{equation}
 1+y_i\leq \sum_k x_{ik}
 \leq 1+(m_{\max}-1)y_i.
 \label{eq:replicas}
\end{equation}
The implementation uses $m_{\max}=2$, so an active task has exactly two
copies. For each fault domain $g$, construction enforces
\begin{equation}
 \sum_{k:g_k=g}x_{ik}\leq 1.
 \label{eq:domains}
\end{equation}
This is a replica-separation requirement. It does not imply that the default
verification set includes whole-domain outages.

The capacity model constrains the cumulative footprint of deployed task
copies:
\begin{align}
 \sum_i\sigma_i x_{ik}&\leq\overline C_k, & k&\in\nodeset,\\
 \sum_i\mu_i x_{ik}&\leq\overline M_k, & k&\in\nodeset.
 \label{eq:capacity}
\end{align}
These inequalities reserve deployment resources across tasks; they are not
instantaneous resource constraints that release capacity after every task
completion. This assumption is maintained during recovery placement, which
starts from the original deployment footprint and adds reexecution demand.

\subsection{Nominal Schedule Evaluation}
The evaluator traverses tasks in a deterministic topological order and
serializes their copies on each node. With $a_k$ denoting node availability,
the data-ready time of copy $(i,k)$ is
\begin{equation}
 b_{ik}=\max_{j\in\Pred(i)}\min_{q\in\copies{j}}
       \{f_{jq}+\tau_{qk}(d_{ji})\},
 \label{eq:data_ready}
\end{equation}
where an empty predecessor maximum is zero. Start and finish times follow
\begin{equation}
 s_{ik}=\max\{a_k,b_{ik}\},\quad f_{ik}=s_{ik}+p_{ik},\quad a_k\leftarrow f_{ik}.
 \label{eq:nominal_schedule}
\end{equation}
Task completion uses the earliest copy,
\begin{equation}
 F_i=\min_{k\in\copies{i}}f_{ik},\qquad \cbase=\max_i F_i.
 \label{eq:nominal_makespan}
\end{equation}
Copies contribute their full modeled compute and communication charges even
when another copy completes earlier. Task-order optimization and idle-gap
insertion are not decision variables of the implemented evaluator.

\subsection{Controllable Fault Scenarios}
The default scenario set depends on the deployment:
\begin{equation}
 \faultset=\Omega^{\mathrm{copy}}(\chi)\cup
 \Omega^{\mathrm{node}}(\chi)\cup\Omega^{\mathrm{link}}(\chi).
 \label{eq:fault_set}
\end{equation}
It contains one copy-failure scenario for every deployed copy and one
node-failure scenario for every occupied node. For each workflow edge whose
effective producer and consumer execute on different nodes, it also includes
an associated transfer-failure scenario. Effective nodes are determined by
the earliest nominal copies in~\eqref{eq:nominal_makespan}; the default set
does not enumerate every possible replica-to-replica transfer failure.

Scenario times are prescribed from the nominal schedule. A copy fails at its
planned finish time; a node fails at its first planned task start; a transfer
fails at its planned completion. Only one such event is injected in a
scenario. Correlated domain failures can be evaluated separately but are not
part of~\eqref{eq:fault_set}. Thus, the term ``worst case'' below means the
largest delay within this explicit, finite scenario set. It does not quantify
over arbitrary failure times, all combinations of failures, or all physical
communication events.

\subsection{Deterministic Recovery Policy}
For each $\omega\in\faultset$, the simulator performs a fault-aware forward
replay in the same topological order. A directly affected active task with
healthy copies completes according to
\begin{equation}
 F_i^{\omega}=\min_{k\in\mathcal H_i^{\omega}}
 \max\{\widetilde f_{ik}^{\omega},
       t_{\omega}+\delta_d+\delta_s+\delta_{\mathrm{sync}}\},
 \label{eq:active_recovery}
\end{equation}
where $\mathcal H_i^{\omega}$ contains surviving copies and
$\widetilde f_{ik}^{\omega}$ is the copy finish time in the replay.
The constants represent detection, switching, and synchronization delay.

When no usable copy remains, a healthy node $k'$ is selected for reexecution.
It must satisfy residual footprint constraints and the policy's domain
restrictions. If no such node exists, the scenario makespan is infinite.
Otherwise,
\begin{equation}
\begin{split}
 F_i^{\omega}={}&\max\{b_{ik'}^{\omega},a_{k'}^{\omega},
                 t_{\omega}+\delta_d+\delta_r\}\\
               &+q_{k'}+p_{ik'},
\end{split}
\label{eq:passive_recovery}
\end{equation}
where $\delta_r$ is replanning delay and $q_{k'}$ is a configured recovery
queue estimate. Node availability $a_{k'}^{\omega}$ includes both healthy
tasks already replayed and earlier reexecutions; omitting it would allow
overlapping execution on a serialized node. Recovery input routing uses the
effective predecessor retained by the replay. This differs from nominal
input selection over all predecessor copies in~\eqref{eq:data_ready}.

Recovery-node selection minimizes a configured heuristic:
\begin{equation}
\begin{split}
 Q_i(k)={}&\theta_1q_k+
 \theta_2\sum_{j\in\Pred(i)}\tau_{\kappa_j k}(d_{ji})
 +\theta_3p_{ik}\\
 &+\theta_4h_i(1-e^{-\lambda_kp_{ik}})
 +\theta_5\rho_kp_{ik},
\end{split}
\label{eq:recovery_selection}
\end{equation}
where $\kappa_j$ is the recovered effective node of predecessor $j$, and
$h_i\in[0,1]$ is structural criticality. Link failures introduce a configured
wait and retransmission delay before affected input becomes available.
This deterministic policy is part of the evaluator, not an optimal recovery
schedule computed independently for each deployment.

\subsection{Recovery SLA and Effective Budget}
Let $C^{\omega}$ be the replay makespan. Define
\begin{equation}
 R_{\omega}=\pos{C^{\omega}-\cbase},\qquad
 \rmax=\max_{\omega\in\faultset}R_{\omega}.
 \label{eq:recovery_delay}
\end{equation}
A complete, capacity-feasible deployment must satisfy
\begin{equation}
 \cbase\leq D,\qquad
 \rmax\leq\reff=\min\{R_{\mathrm{SLA}},D-\cbase\}.
 \label{eq:sla}
\end{equation}
Its recovery safety margin is $M(\chi)=\reff-\rmax$.
Figure~\ref{fig:budget} illustrates why nominal slack and recovery tolerance
must be considered together.

\begin{proposition}
For a complete, capacity-feasible deployment and the specified scenario set,
\eqref{eq:sla} holds if and only if $\cbase\leq D$, every
$C^{\omega}\leq D$, and every $R_{\omega}\leq R_{\mathrm{SLA}}$.
\end{proposition}
\begin{IEEEproof}
If~\eqref{eq:sla} holds, then
$C^{\omega}\leq\cbase+\rmax\leq D$ and
$R_{\omega}\leq\rmax\leq R_{\mathrm{SLA}}$.
Conversely, $C^{\omega}\leq D$ and $\cbase\leq D$ imply
$R_{\omega}\leq D-\cbase$. Combining this with the recovery SLA and taking
the maximum over scenarios gives~\eqref{eq:sla}.
\end{IEEEproof}

\begin{figure}[!t]
 \centering
 \begin{tikzpicture}[x=0.68cm,y=0.63cm,>=Stealth,font=\footnotesize]
 \fill[gray!16] (0,0)--(5.4,0)--(3.4,2.0)--(0,2.0)--cycle;
 \draw[->] (0,0)--(6.0,0) node[below right] {$\cbase$};
 \draw[->] (0,0)--(0,3.2) node[above] {$\rmax$};
 \draw[thick] (0,2.0)--(3.4,2.0)--(5.4,0);
 \draw[dashed] (3.4,0)--(3.4,2.0);
 \node[left] at (0,2.0) {$R_{\mathrm{SLA}}$};
 \node[below] at (5.4,0) {$D$};
 \node[below] at (3.4,0) {$D-R_{\mathrm{SLA}}$};
 \node at (1.7,0.85) {Feasible};
 \node[rotate=-43,above] at (4.45,1.05) {$D-\cbase$};
 \node[align=center] at (3.0,2.85) {Nominal makespan and additional delay\\are measured in seconds};
\end{tikzpicture}

 \caption{Conceptual feasible timing region. The effective recovery limit is
 the smaller of the declared window and the remaining deadline slack. This
 diagram illustrates the model; it is not an experimental plot.}
 \label{fig:budget}
\end{figure}

\subsection{Cost Objective}
The monetary component of nominal deployment is
\begin{align}
 J_{\mathrm{comp}}&=\sum_i\sum_{k\in\copies{i}}\rho_kp_{ik},\\
 J_{\mathrm{comm}}&=\gamma_c\sum_{(i,j)\in\edgeset}d_{ij}
      \sum_{k\in\copies{i}}\sum_{\ell\in\copies{j}}\mathbf1[k\ne\ell],\\
 J_{\mathrm{res}}&=\gamma_{\mathrm{res}}\sum_{i:y_i=0}\sigma_i.
 \label{eq:cost_terms}
\end{align}
Communication billing counts all provisioned cross-node replica pairs, even
though input readiness uses the fastest nominal predecessor copy. This is an
explicit billing convention. Replica compute is already included in
$J_{\mathrm{comp}}$ and is not added again as a separate backup charge.

Let $J_0=J_{\mathrm{comp}}+J_{\mathrm{comm}}+J_{\mathrm{res}}$. The reported
objective adds a recovery-delay proxy:
\begin{equation}
 J(\chi)=J_0(\chi)+\gamma_{\omega}\overline R(\chi),\qquad
 \overline R=\frac{1}{|\faultset|}\sum_{\omega\in\faultset}R_{\omega}.
 \label{eq:objective}
\end{equation}
For feasible deployments, all terms of this average are finite. It is a
uniform scenario-average delay with a configured conversion weight, not a
probability-weighted expected financial loss. We use it for feasible-solution
comparison and do not interpret the simulator's finite-value proxy on
unrecoverable deployments as an expectation.

RLC-HFTSP minimizes~\eqref{eq:objective} over placement and protection modes
subject to~\eqref{eq:replicas}--\eqref{eq:capacity}, nominal dependency
evaluation, and~\eqref{eq:sla}. The recovery policy is fixed. The problem is
a combinatorial deployment search; no claim of global optimality or an
unproved complexity reduction is required by the proposed heuristic.
